\section{Discussion}
\label{sec:discussion}

\subsection{Principal findings}
On DROZY, point estimates of the error were lower with each added source (MAE 1.73 for the prior alone, 1.38 with a check-in and passive signals), but the intervals overlap, the prediction sets differ, the sample is small, and a protocol-aware naive baseline did as well. All of the gain arose at the night test, where the prior was most biased, and part of it is likely calibration transfer within one laboratory. On ds004902, a group-mean offset predicted ratings better than a person's own offset from another session, and the KSS mapping refitted by Ingre et al.~\cite{ingre2014validating} reduced the prior's error by 0.22 points, although DROZY's error lies in the prior's time course rather than its mapping. The website's camera pipeline, run on DROZY's near-infrared video, correlated with EOG blink rate but undercounted blinks by about a third; its blink-duration association with sleepiness was confounded with frame rate; and its PERCLOS failed, so the model now ignores PERCLOS in that case. In simulation, check-ins carried most of the information, and passive signals alone did not improve on the prior. Estimating per-signal anchors gave a small gain (0.03 KSS points over check-ins only, 0.04 over a filter that allows for persistent biases without estimating them), which held when the model's assumptions were broken. In a simulation of the planned 10-day study, replayed as one continuous log, passive signals improved the prediction of observed ratings by 0.01 points, and a mean-reverting offset by a further 0.02 to 0.04; when each day was replayed separately, so that the anchors were relearnt every morning, the passive signals made predictions slightly worse. The value of the anchors therefore depends on carrying them across days. The evidence supports check-ins as the main source of information and anchored passive signals as a small addition whose value in real use is unproven.

\subsection{Comparison with other work}
Table~\ref{tab:comparison} lists published results for related problems. No study we found reports KSS regression error on DROZY or ds004902; published DROZY results are binary classifications~\cite{maior2020realtime}, which cannot be compared with an MAE, so our numbers are a first reference point on these datasets rather than evidence that BlinkTwin is more accurate than other systems. On the KSS scale, our errors are in the range expected for the Three-Process Model: Ingre et al.\ reported an in-sample residual SD of 1.46 in field data~\cite{ingre2014validating}, and Knock et al.\ predicted 48--60\% of shift workers' ratings within one point~\cite{knock2021prediction}, against 36\% for our prior on ds004902 (45\% with the refitted mapping). The closest designs predict PVT performance rather than self-rated sleepiness. Subramaniyan et al.\ fed face-video features into a personalised model and reached an error close to the PVT's own within-person variability~\cite{subramaniyan2025personalized}, the 2B-Alert app needed about 12 PVTs to reach its best accuracy~\cite{reifman2019twobalert}, and Ahlstrom et al.\ improved a sleep/wake predictor's forecast of high KSS with an eye-movement test~\cite{ahlstrom2013fit}. That personal learning needs many observations is consistent with our finding that one or two check-ins help little. Accuracy on real drowsiness from consumer cameras is modest: on UTA-RLDD the best video model reached 65.2\% on three classes and human judges 57.8\%~\cite{ghoddoosian2019realistic}.

\subsection{Subjective and objective sleepiness}
Anchoring makes the person's ratings the reference, which is right for an estimate of how sleepy a person feels but can hide objective decline~\cite{vandongen2003cumulative,leproult2003individual}. A single reaction test read on the population scale cannot reveal that decline, because people's baselines differ by about 2.6 KSS-equivalents and device latency shifts the level. The application therefore flags a possible dissociation only from within-person change in the reaction-time anchor, and presents the check as experimental. Its false-alarm rate, detection delay and effect on behaviour are untested; a chronic-restriction scenario, in which ratings level off over days while performance declines, is the right test.

\subsection{How accuracy could improve}
In order of the evidence we have: (i)~a mean-reverting offset, which lowered simulated error without loss of coverage in the one-day study, improved rating prediction most in the 10-day power simulation, and matches the weak persistence in ds004902; (ii)~a personal state beyond an additive offset, such as a gain on homeostatic pressure or a circadian phase shift, which trait-like vulnerability~\cite{vandongen2004systematic} predicts and which two sessions at different sleep pressure can identify; (iii)~the field-refitted KSS mapping, supported by ds004902 but not by DROZY; (iv)~forward models refitted on the website's own camera features, with the closed-eye level normalised to the person, the blink-duration noise inflation matched to that pipeline, long closures as a further signal, and the frame rate, which the application now stores with every minute, used to exclude or model low-rate minutes; (v)~a lapse- or speed-based PVT model fitted on the PVT-B itself; and (vi)~check-ins scheduled where they are most informative, which the application now suggests.

\subsection{Limitations}
The real-data evaluation rests on 14 people in DROZY and 33 in ds004902, with near-infrared rather than RGB video, a 10-min rather than 3-min PVT, assumed wake times in DROZY, frame rate confounded with sleep deprivation, and landmark rather than MediaPipe features in the replay. The forward-model slopes are fitted on noisy, integer ratings and are therefore attenuated~\cite{carroll2006measurement}, and slope heterogeneity is treated as independent noise although it is a persistent person-level error. The EOG reference was an automatic detector without manual annotation, and event-level blink and long-closure accuracy was not measured. Browser reaction times include device latency that we have not yet measured. The KSS is treated as a continuous Gaussian measurement although it is an integer scale with a ceiling, and the model's interval describes the latent state, not a future rating, unless rating noise is added. The prior uses local clock time without phase adaptation, so travel, shift work and daylight-saving changes produce time-varying errors that the offset cannot absorb. The time-on-task term, the offset's dynamics and the alert thresholds are assumptions, and the forward-pitch sign of the head-drop alert has not been verified on devices. The simulation, although it includes misspecified scenarios and a fair comparator, is built from the model's own structure, and the 24-h scenario lies outside the intended use; on an ordinary workday sleepiness of KSS~7 or more was rare. The privacy check used a test pattern without a face. Workload was not evaluated. BlinkTwin is not a medical device.

\subsection{Ethics and privacy}
Webcam monitoring at work raises concerns about surveillance. BlinkTwin is designed for self-monitoring: frames never leave the page, a content security policy blocks all connections other than to the site and the model host, only one-minute numeric summaries are stored in the person's browser (unencrypted, which matters on shared computers), and the person can export or delete them. Any deployment by an employer would need informed consent, transparency about what is measured, and a guarantee that the data are not used against the person.

\section{Conclusion}
\label{sec:conclusion}
BlinkTwin estimates how sleepy a person feels by joining the published Three-Process Model with occasional KSS check-ins and with reaction-time and camera signals that each have their own personal anchor, entirely in a browser and with video kept on the device. On public data, check-ins and the prior's time course mattered most, the camera pipeline measured blinks but not PERCLOS on near-infrared video, and the passive signals' benefit was small and confined to simulation. Validation of the camera pipeline on ordinary webcams and the planned prospective study are needed before any claim about daily use.
